% page 272 of the facsimile (image p-271 of the scan)
% block 160.0 x 250.6 mm ; left margin 8.0 mm
\pgc{-12.700mm}{0.000mm}{
\l{}
\l{}
\l{\cel{4}264}
\l{}
\l{}
\l{\sou{karamelo}. Sukro fuzita e brunigita per la efiko da fairo.}
\l{\cel{3}- DEFRS.}
\l{}
\l{\sou{karasino}. (\sou{zool.)}\cel{2}Fisho simila kam karpo. - L.\cel{1}\sou{carassius}}
\l{\cel{3}\sou{vulgaris}.}
\l{}
\l{\sou{karato.}\cel{2}Pondero-unajo, uzata ankore da la juvelisti, por la}
\l{\cel{3}diamanti, la lapidi, la perli, e qua equivalas cirkume}
\l{\cel{3}20 centigrami. - DEFIRS.}
\l{}
\l{\sou{karavano}. - I. En Avan-Azia, en Arabia, e c., trupo de komer-}
\l{\cel{3}cisti, de voyajeri, de pilgrimanti, asemblita por trans-}
\l{\cel{3}irar plu sekure la dezerti. - II.(\sou{metaf.)}\cel{1}Trupo de personi}
\l{\cel{3}qui iras ensemble. - DEFIRS.}
\l{}
\l{\sou{karavanserayo}.\cel{2}En la Oriento-landi : granda edifico en qua la}
\l{\cel{3}voyajeri povas shirmar e su ipsa e lia charjo-bestii. DEFIRS}
\l{}
\l{\sou{karavelo}.\cel{2}I. (\sou{olim)}\cel{1}. Che la Turki, grosa milito-navo. -}
\l{\cel{3}II. Mikra navo kun segli Latina, uzita precipue da la Por-}
\l{\cel{3}tugalani. - III. Navo uzata por peskar la haringi. - F.}
\l{}
\l{\sou{karbo}. (\sou{kemio}).Korpo simpla, metaloida, qua existas plura-}
\l{\cel{3}forme : karbono, grafito, diamanto, e c. -\cel{1}\sou{Simbolo kemiala}:}
\l{\cel{3}C. - DEFIS.}
\l{}
\l{\sou{karbono}. Materio nigra, en qua dominacas la korpo simpla}
\l{\cel{3}\textquotedbl{}karbo\textquotedbl{}, e quan onu extraktas (ordinare) de subsulo per}
\l{\cel{3}minar. -EFIS.}
\l{}
\l{\sou{karbonaro}.\cel{2}Membro di societo revolucionista sekreta, fondita}
\l{\cel{3}en Italia, komence di la l9esma yarcento. - FI.}
\l{}
\l{\sou{karbonario}.\cel{2}(\sou{zool.)}\cel{2}Fisho ek la genero \textquotedbl{}gado\textquotedbl{}. - L.\cel{1}\sou{gadus}}
\l{\cel{3}\sou{carbonarius}..}
\l{}
\l{\sou{karborundo}. (\sou{kemio}). Siliko-karburo, preparita per submisar,}
\l{\cel{3}a la kaloro de elektroforno, mixuro ek karbono e silicio:}
\l{\cel{3}uzata por la fabriko di grindo-roti, e c.}
\l{}
\l{\sou{karbunklo}.\cel{2}I. (\sou{patol.})\cel{2}Morbo qua desorganizas e nigrigas la}
\l{\cel{3}tisui atakita. - II. Varietato di granato reda, di qua la}
\l{\cel{3}brilo esas tre granda. - DEFIRS.}
\l{}
\l{\sou{karburar}. (\sou{netrans.}) (\sou{tekn.)}\cel{2}Preparar la mixuro gasa,explo-}
\l{\cel{3}zala quan on destinas a la alimento di la explozo-motori. EF}
\l{}
\l{\sou{karburatoro}. (\sou{tekn.)}\cel{2}Aparato qua produktas ta mixuro ek aero}
\l{\cel{3}e vaporo de hidrokarburo o kombusto-gaso, quan onu uzas}
\l{\cel{3}sive por lumizo, sive por alimentar explozo-motori. - EF.}
\l{}
\l{karcero. Loko ube la arestito retenesas, malgre lua volo e,}
\l{\cel{3}tale, privacesas de lua libereso. - DEFIRS.}
}
